% !TeX program = pdflatex
%
% ux-paper.tex
%
% Build:  pdflatex ux-paper.tex   (run 3x: ToC, then cross-refs)
%

\documentclass[11pt,a4paper]{article}

\usepackage[T1]{fontenc}
\usepackage[utf8]{inputenc}
\usepackage[margin=2.4cm]{geometry}
\usepackage{amsmath}
\usepackage{amssymb}
\usepackage{booktabs}
\usepackage{longtable}
\usepackage{array}
\usepackage{xcolor}
\usepackage{enumitem}
\usepackage{fancyhdr}
\makeatletter\def\input@path{{../}{../../design/brand/}}\makeatother
\usepackage{isocline}   % docs/research/isocline.sty — Isocline Research template
\usepackage[hidelinks]{hyperref}

% Legacy markup, now neutral: Isocline reports observations, not verdicts.
\newcommand{\bad}[1]{\obs{#1}}
\newcommand{\good}[1]{#1}
\newcommand{\code}[1]{\texttt{\small #1}}

\isoclineshorttitle{User Experience Design}

\title{\textbf{Progressive Clarity with Honest Constraint:}\\[0.3em]
\large A Design Framework for Hormone Replacement Therapy\\
Tracking in the Isocline Mobile Application}
\author{Starla Cortes}
\date{August 2026}

\begin{document}
\maketitle

\begin{abstract}
\noindent
We present the user experience architecture of Isocline, a mobile application
that predicts serum hormone concentration over time from a user's logged
dose history for gender-affirming hormone replacement therapy (HRT). The
application translates a pharmacokinetic (PK) model spanning ten
medication/route/form combinations into a consumer interface designed for
daily use by people without medical training.

The UX framework rests on three principles: \textbf{progressive
disclosure}, which reveals complexity only as the user demonstrates
readiness---from an 11-step onboarding wizard through to a 4-tab daily
experience with tenure-tiered content that unlocks over months and years
of use; \textbf{honest presentation}, which reports hormone levels as
population estimates rather than facts, gates comparative judgments
behind user-configured therapeutic targets, and phrases safety flags as
clinical associations rather than directives; and \textbf{per-user
personalization}, which adapts the interface to the individual through
onboarding choices, blood test calibration, body-type selection, and
injection-zone preferences.

We describe the full user journey from first launch through ongoing
mastery, document the design-token system that encodes clinical semantics
as visual properties, validate the architecture against eight persona
journeys that exercise UX edge cases, and provide an honest accounting of
design trade-offs: where the UX is anchored, where it is reasoned but
untested, and where it remains placeholder. The paper is written to be
checkable: every design decision is cited to its source in the codebase,
and every limitation is named.

\textbf{Keywords:} user experience, mobile health, progressive
disclosure, hormone replacement therapy, design system, onboarding,
personalization, pharmacokinetic visualization
\end{abstract}

\tableofcontents
\clearpage

% ======================================================================
\section{Introduction}
\label{sec:intro}

Gender-affirming hormone replacement therapy (HRT) involves the
administration of exogenous hormones---estradiol and/or
testosterone---to align an individual's serum hormone levels with their
gender identity. Clinical management relies on periodic blood draws to
measure serum hormone concentration, titrate dosage, and confirm
therapeutic levels~\cite{hembree2017}. Between draws, patients have limited
tools to anticipate where their levels are heading, when a peak or trough
is expected, or how a missed dose or regimen change will affect the
trajectory.

Isocline addresses this gap with a pharmacokinetic model that predicts serum
hormone concentration over time from the user's logged dose history. The
model itself is described in a companion paper~\cite{xtpk}; this paper
concerns the \emph{interface} that delivers that model to its users. The
challenge is non-trivial: the underlying mathematics are undergraduate-level
differential equations, the clinical context is high-stakes, and the
audience is anyone on HRT---not just those comfortable with exponential
decay curves.

\subsection{The UX Problem}
\label{sec:ux_problem}

People on HRT must manage a constellation of tasks: tracking injection
schedules across dosing intervals that range from daily to quarterly,
rotating injection sites across body zones to prevent tissue damage,
interpreting hormone level estimates in the context of therapeutic
ranges, correlating subjective experience (mood, symptoms) with objective
data (blood test results), and anticipating physical changes over months
and years. None of these tasks require medical training, but all of them
benefit from clinical context.

Most available tools fall into one of two categories. General-purpose
medication reminders reduce the problem to a binary ``did you take it''
checkbox, discarding the pharmacokinetic information that makes sense of
\emph{when} a level is measured and \emph{where} in the cycle the user
sits. Clinical reference materials, at the other extreme, present PK
curves and population tables without translation---accurate but
illegible to someone who has never seen a biexponential.

Isocline occupies the space between: it computes a full pharmacokinetic
prediction but presents it as a single number, a single curve, and a
single status. The complexity is present but gated; detail is available
on demand but never imposed.

\subsection{The Design Response}
\label{sec:design_response}

Isocline's interface is organized around three principles that operate at every
level of the application, from the 11-step onboarding wizard to the
ambient constellation visualizations on the Time tab.

\textbf{Progressive disclosure.} The app never shows more than the user
needs at their current stage. Onboarding introduces one concept per step.
The home screen foregrounds exactly one number. The Time tab's content is
gated by tenure tier---a user in month~1 does not see year-3 surveillance
content. Detail lives behind taps, and taps are only offered when the
detail is relevant.

\textbf{Honest presentation.} The PK model produces population estimates,
not individual predictions. The interface reflects this at every level:
level readouts are labeled ``Estimated,'' comparative judgments (``below
target'') only appear against user-configured therapeutic targets, and
safety flags state clinical associations (``Estradiol above
1,000\,pg/mL is associated with increased thromboembolic risk'') rather
than directives (``You should lower your dose''). The app reports what
the model says; it does not advise.

\textbf{Per-user personalization.} Between-subject variability in hormone
metabolism is substantial---published studies report coefficients of
variation of 0.58 on $V_d/F$ for injectable testosterone~\cite{ichihara2020},
corresponding to an $\sim$8-fold individual range. The interface
addresses this through a chain of personalization: medication type and
ester at onboarding, a per-user calibration factor $\alpha$ derived from
the user's own blood tests, body-type selection for the 3D injection-site
model, and user-configurable therapeutic target bands. No single setting
carries the full weight of individualization; each layer tightens the
fit.

\subsection{Contributions}
\label{sec:contributions}

\begin{enumerate}[nosep,leftmargin=1.4em]
\item A complete walkthrough of the UX architecture, from 11-step
  onboarding through the 4-tab daily experience, with design rationale
  for each structural decision.
\item A design-token system that encodes clinical semantics (level status,
  safety flags, cycle position, tenure tier) as visual properties without
  presenting medical advice, implemented in a single presentation-rules
  file (\code{HormoneLevelPresentation.swift}) that separates semantics
  from display.
\item A progressive disclosure framework that operates on three
  timescales simultaneously: the onboarding sequence (minutes), the
  per-tab depth gradient (seconds per interaction), and the tenure-tier
  system (months to years).
\item A per-user calibration interface that translates a blood test result
  into a single multiplicative factor $\alpha$, absorbing the dominant
  source of individual-level prediction error without requiring the user
  to understand population pharmacokinetics.
\item An honest accounting of UX limitations---what is designed well,
  what is reasoned but untested, and what remains unaddressed---in the
  tradition of the companion PK paper's explicit marking of calibrated
  versus uncalibrated parameters.
\end{enumerate}

% ======================================================================
\section{Design Principles}
\label{sec:principles}

This section describes the five principles that govern every design
decision in Isocline. They are the UX analog of the PK model's governing
equation (Equation~1 in~\cite{xtpk}): every screen, every animation, and
every string of user-facing text is answerable to at least one of them.

\subsection{State Facts, Don't Advise}
\label{sec:state_facts}

The central tension in any health-tracking application is the line
between \emph{informing} and \emph{advising}. Isocline's position is that the
line exists, that crossing it requires clinical authority the app does
not possess, and that the interface must be designed to stay on the
correct side of it.

This principle is implemented in \code{HormoneLevelPresentation.swift},
the single file through which every hormone level estimate becomes
user-facing text. The file exports classification enums
(\code{LevelStatus}, \code{SafetyFlag}, \code{ModelDeviation},
\code{RangeSource}) and a set of formatting functions. Its design
follows three rules:

\begin{enumerate}[nosep,leftmargin=1.4em]
\item \textbf{Classifications are values, not display strings.} The
  \code{LevelStatus} enum has cases \code{criticallyLow},
  \code{veryLow}, \code{slightlyBelow}, \code{inRange},
  \code{slightlyAbove}, and \code{high}---but these are model facts, not
  verdicts. Views build their own text from these values, which keeps
  semantics and presentation (and therefore localization) separate. The
  app never ships a hardcoded English string like ``Your levels are
  low'' from the model layer.

\item \textbf{Comparisons are gated on user-set bounds.} The field
  \code{RangeSource} distinguishes \code{userSet} bounds (the user has
  configured therapeutic targets in Settings) from \code{modelDerived}
  bounds (population defaults). Comparative copy---``within your set
  trough range,'' ``below your set target''---is only rendered when the
  source is \code{userSet}. If the user has not configured targets, the
  app reports the estimate without judgment. The model may know that
  80\,pg/mL is below the typical estradiol trough, but the app does not
  tell the user they are ``below target'' against a target they never
  set.

\item \textbf{Safety flags state associations, not instructions.} The
  \code{SafetyFlag} enum identifies crossings of documented clinical
  thresholds: \code{estrogenCeiling} (E2 ${}> 1000$\,pg/mL),
  \code{testosteroneSupraphysiologic} (T ${}> 1500$\,ng/dL),
  \code{subTherapeuticTrough}, and
  \code{spironolactoneClearance}. The user-facing text for each flag
  states the association (``Estradiol above 1,000\,pg/mL is associated
  with increased thromboembolic risk'') without prescribing action. The
  app's position is that stating the association is informing; saying
  ``you should lower your dose'' is advising, and the app does not do
  it.
\end{enumerate}

This principle has a cost: it makes the interface less directive than
some users might prefer. A user who wants the app to say ``take your dose
now'' or ``your levels are too high---skip today'' will find Isocline
reticent. That reticence is intentional. The companion PK
paper~\cite{xtpk} notes that the model's predictions ``are population
estimates that should not be used to make dosing decisions without
laboratory confirmation''; the interface honors that constraint by
refusing to translate population estimates into individual
instructions.

\subsection{Progressive Disclosure}
\label{sec:progressive}

Isocline practices progressive disclosure on three timescales simultaneously.

\textbf{Onboarding (minutes).} The 11-step onboarding wizard
(Section~\ref{sec:onboarding}) introduces exactly one concept per screen:
first the hormone type, then the experience level, then the
administration method, and only then the ester---a term many users
encounter for the first time in the app. Each step builds on the
selections made in prior steps: the ester picker only appears after
``Injectable'' is chosen; the concentration picker is pre-populated with
common values for the selected ester; the injection-site step is skipped
entirely for oral routes. The user is never asked a question whose
context they have not already been given.

\textbf{Daily use (seconds per interaction).} The home screen foregrounds
exactly one number: the current estimated hormone level, displayed in
large type with units and a status descriptor. The full PK curve is
present as a background layer---visible but ambient, not demanding
attention. The glass card stack below it contains four rows, each a
single piece of information: insight text, today's mood, next dose
countdown, and action buttons. Tapping any row opens detail; until
tapped, the row is a summary. The design principle is that the home
screen should be legible in under two seconds and actionable in under
five.

\textbf{Tenure (months to years).} The Time tab's content is gated by
\code{TenureTier}, a 7-level progression from \code{beginning}
(0--2\,months) through \code{veteran} (120+\,months). Each tier unlocks
specific content about expected changes: a user in month~1 sees
information about initial physical responses; a user at year~3 sees
content about long-term health surveillance. This gating is not
gamification---the tiers are not displayed as achievements to chase---but
a content-management strategy: a user in their first month does not need
to read about year-3 mammogram schedules, and presenting that information
prematurely would be noise, not signal.

The tiers also control the visual complexity of the Time tab's
constellation point-clouds (Section~\ref{sec:time_tab}). The Mind and
Body clouds grow denser and more articulated as the user's tenure
increases, reflecting the accumulation of data and experience. A user
with 3~days of dose history sees a sparse field; a user with 3~years sees
a rich, structured visualization. The change is gradual enough to be
subliminal---most users will not consciously notice the cloud
density---but the effect is that the Time tab deepens with use rather
than stagnating.

\subsection{Honest Uncertainty}
\label{sec:uncertainty}

The PK model's predictions are population estimates. The companion
paper~\cite{xtpk} quantifies this: out-of-sample validation errors of
1--15\% for estradiol and testosterone single-dose and steady-state
checks, with specific structural limitations (TU peak/trough ratio,
transdermal intraday ripple, oral testosterone peak fluctuation) that
bound what the model can honestly claim.

The interface reflects this uncertainty in five ways:

\begin{enumerate}[nosep,leftmargin=1.4em]
\item The level overlay on the home screen reads ``Estimated'' before the
  numeric value---a one-word qualifier that appears on every readout,
  establishing that the number is a model output, not a measurement.
\item The calibration factor $\alpha$ is displayed in Settings after a
  blood test is entered, making the personalization explicit rather than
  implicit. The user can see that the model has been adjusted, and by how
  much.
\item The PK model's known limitations---TU peak/trough ratio,
  transdermal intraday ripple, TP's uncalibrated $t_{\max}$---are
  \emph{not surfaced in the consumer interface}. This is itself a design
  decision: the app does not burden the user with caveats about every
  prediction. The limitations are documented in the companion paper and
  in code comments; they constrain what the app claims, but they are not
  presented as consumer-facing disclaimers. This choice is discussed
  further in Section~\ref{sec:limitation_surfacing}.
\item The therapeutic target bands on the graph are only rendered when the
  user has configured them. A population-default band could be drawn
  automatically---the model knows the typical ranges---but doing so would
  present a clinical reference that the user did not request, blurring
  the line between information and advice.
\item The graph's visible time window and sampling resolution are derived
  from the ester's pharmacokinetic properties
  (Equations~10--11 in~\cite{xtpk}), not hardcoded. This means the curve
  always shows a clinically appropriate window for the specific
  medication, but it also means the window changes when the medication
  changes---a subtle consistency trade-off that privileges accuracy over
  familiarity.
\end{enumerate}

\subsection{Per-User Personalization}
\label{sec:personalization}

Between-subject variability in hormone metabolism is substantial, and the
app addresses it through a chain of personalization layers, each narrower
than the last.

\textbf{Onboarding configuration.} The medication type (estrogen vs.\
testosterone), ester, dose, concentration, and frequency are the coarsest
layer. They select which PK constants to use and anchor the population
prediction to the user's reported regimen. These choices are
high-confidence: they come directly from the user's prescription.

\textbf{Blood test calibration.} The calibration factor
$\alpha = C_{\text{measured}} / C_{\text{predicted}}$
(Section~3.4 of~\cite{xtpk}) is a single scalar multiplier on the entire
model output, derived from the user's own blood test results. The
interface for entering a blood test is minimal: date, measured level,
hormone type. The app computes $\alpha$ silently and applies it to all
subsequent predictions. The user does not need to understand the factor
to benefit from it; the curve simply becomes more accurate after a blood
test is recorded.

\textbf{Body type.} The 3D body model used for injection-site tracking
offers three silhouettes (neutral, feminine, masculine). This is a visual
personalization that affects only the body model's mesh; it does not
alter the PK predictions. Its presence in onboarding---step~8 of
11---signals that the app acknowledges anatomical diversity without
making it medically consequential.

\textbf{Therapeutic targets.} Users can configure trough minimum, trough
maximum, and peak ceiling values in Settings. These numbers draw the
reference bands on the home screen graph and gate the comparative text in
the level overlay. They are the only user-configurable clinical
parameters in the app, and they are explicitly optional: if not set, the
app reports estimates without comparison.

\textbf{Injection zones.} Users can mark body zones as ``preferred'' or
``avoid'' on the 3D model. The app uses these preferences to suggest
rotation patterns, but---critically---it does not enforce them. The user
can always tap anywhere on the body; the zone preferences are advisory,
not restrictive.

\subsection{Accessibility as First Principle}
\label{sec:accessibility}

Isocline's motion and typography systems are designed with accessibility
constraints as defaults, not as overrides to be applied later.

\textbf{Reduce Motion.} Every animation in the app is gated on
\code{@Environment(\.accessibilityReduceMotion)}. The
\code{EntranceModifier} (Section~\ref{sec:motion}) checks Reduce Motion
on appear; when enabled, views appear immediately at their final position
with no animation. The confetti burst on dose recording is suppressed.
The constellation point-clouds on the Time tab freeze into static
arrangements. The spring-driven button press animation degrades to an
instant opacity change. The design goal is that no information is ever
conveyed \emph{only} through motion; motion is decoration, and
decoration must be removable.

\textbf{Dynamic Type.} The app uses role-based font tokens rather than
raw \code{.system(size:)} literals. The tokens map to Apple's built-in
Dynamic Type text styles (\code{.title2}, \code{.body}, \code{.caption},
etc.), which scale automatically with the user's system font size
preference. Custom tokens (\code{Theme.Fonts.statValue},
\code{Theme.Fonts.hero}, \code{Theme.Fonts.readout}) are built on top of
these styles with weight and design adjustments, preserving the scaling
behavior. The \code{@ScaledMetric} property wrapper is used for
layout-sensitive values like horizontal padding, so that spacing grows
proportionally with text.

\textbf{Dark-mode-first design.} Every color in the design system is
defined via \code{Color(lightHex:, darkHex:)}, with the dark variant
treated as the primary and the light variant as the adaptation. The
brand gradient, graph background, glass surfaces, and semantic colors all
have explicit light-mode counterparts. The app does not rely on
Apple's automatic color inversion; every color pair is authored.

\textbf{VoiceOver.} The app uses standard SwiftUI accessibility modifiers
(\code{.accessibilityLabel}, \code{.accessibilityHint},
\code{.accessibilityValue}) on interactive elements. The 3D body model
presents an accessibility description of selected zones. The graph,
however, is not fully navigable by VoiceOver---a limitation discussed in
Section~\ref{sec:unaddressed}.

% ======================================================================
\section{The User Journey}
\label{sec:journey}

This section traces the complete user experience from first launch
through ongoing mastery. Each subsection describes what the user sees and
does, followed by the design rationale.

\subsection{Onboarding: 11 Steps to Configuration}
\label{sec:onboarding}

The onboarding wizard (\code{OnboardingView.swift}, approximately 1,600
lines) is the user's first interaction with Isocline. It is a linear,
non-branching sequence of 11 screens (10 for oral routes, which skip the
injection-site step), each asking exactly one question. A progress bar at
the top shows position within the sequence. Steps fade in and out with a
0.3-second opacity transition; there is no lateral swiping or skipping
ahead.

The wizard is presented as a full-screen cover over the main interface.
Until it completes, the main interface does not exist from the user's
perspective. This is a deliberate commitment: onboarding is not a tutorial
that can be dismissed and returned to later; it is the configuration
mechanism for the PK engine, and the app cannot function without it.

\subsubsection{Step 0: Welcome}

A single screen with the word ``Isocline'' in 64pt light weight and the
subtitle ``Hormone therapy tracking.'' A single button reads ``Begin.''
There is no account creation, no email collection, no permission request.
The user is not asked to trust the app before they have seen what it
does.

\subsubsection{Step 1: Medication Type}

``What hormone are you taking?'' Two cards: ``Estrogen-based HRT''
(estradiol, feminizing) and ``Testosterone-based HRT'' (masculinizing).
Each card includes a one-line plain-language description. Selection
determines which branch of the PK model will be used and which esters,
units, and therapeutic ranges will be shown in subsequent steps.

\subsubsection{Step 2: Experience}

``How long have you been on HRT?'' Two options: ``Just starting'' or
``I've been on it.'' If the latter, year and month steppers appear (0--20
years, 1--11 months). This value seeds the endogenous baseline model
(Section~3.10 of~\cite{xtpk}) and the Time tab's tenure tier. It is the
only onboarding step that has no visual consequence in the immediate
post-onboarding experience; its effect unfolds over months.

\subsubsection{Step 3: Method and Ester}

``How do you take it?'' Three method cards: Injectable, Cream, Oral. If
Injectable is selected, an ester picker appears below with an animated
expand transition. Each ester card shows the ester name, its abbreviation
(EV, EC, TC, TE, TU, TP), its terminal half-life in days, and its typical
time-to-peak. These are the same constants that appear in
Table~2 of~\cite{xtpk}, translated into plain language: ``Peak: $\sim$2
days'' rather than ``$t_{\max} = 2.10$\,d.''

Why show half-life and peak at onboarding? Because the ester determines
the dosing interval, and the dosing interval determines the user's lived
experience. A user choosing between testosterone cypionate (peak
$\sim$2\,days) and testosterone undecanoate (peak $\sim$7\,days) is
making a decision about how frequently they will inject, and the
interface gives them the information to make that decision without leaving
the app to search.

A frequency stepper (``Every $N$ days,'' 1--30) appears below the method
selection. The default is 7\,days for injectables, reflecting the most
common clinical practice.

\subsubsection{Step 4: Concentration}

``What's your vial concentration?'' For each ester, a hardcoded list of
common concentrations is presented as selectable chips: 40\,mg/mL for
estradiol valerate, 200\,mg/mL for testosterone cypionate, 5\,mg/mL for
estradiol cypionate. The most common concentration for the selected ester
is marked ``$\cdot$ most common.'' A ``Custom'' option opens a text field
for arbitrary values.

This step exists because concentration is \emph{not} derivable from the
ester alone. A user on testosterone cypionate might have a 200\,mg/mL
vial or a 100\,mg/mL vial; the difference is a factor of 2 in the
computed volume, and getting it wrong would make the mg/mL conversion
(Section~\ref{sec:dose_step}) silently incorrect. By making concentration
an explicit onboarding step rather than a buried setting, the app ensures
the conversion is correctly anchored from the start.

\subsubsection{Step 5: Dose and Timing}
\label{sec:dose_step}

``What's your dose?'' This step introduces the mg/mL toggle, one of the
most heavily designed interactions in the app. The problem it solves is
specific to injectable medications: prescriptions are written in
milligrams (e.g., ``4\,mg estradiol valerate''), but syringes are marked
in milliliters. The conversion depends on the vial concentration
($\text{mL} = \text{mg} / \text{concentration}$), and getting it wrong
means injecting the wrong volume.

The step presents a segmented control with two options: ``mg'' and
``mL.'' In mg mode, a stepper and text field accept the dose in
milligrams; the computed volume in mL is shown below in smaller type
(``$=$ 0.10\,mL at 40\,mg/mL''). In mL mode, the fields accept volume
directly and compute the mg amount. Changing the concentration picker
updates both values in real time.

A row of quick-dose chips ($0.5$\,mL, $1.0$\,mL, $1.5$\,mL in mL mode;
$2$\,mg, $4$\,mg, $6$\,mg in mg mode) provides one-tap entry for common
doses. The chip values adjust based on the selected medication: estradiol
shows lower mg values; testosterone shows higher ones.

A date picker sets the date of the user's most recent dose. The
\emph{time} of day is not collected at onboarding---the Quick Dose Sheet
(Section~\ref{sec:first_dose}) handles time---but the date is needed to
place the first data point on the PK curve. The picker defaults to
today, which is correct for users who took their dose on the day they
set up the app.

\textbf{Volume warnings.} If the computed injection volume exceeds
3.0\,mL (the typical maximum for a single intramuscular injection), a
warning appears: ``This volume is unusually large for a single injection.
Check your dose or concentration.'' This is a safety net, not a
blocker---the user can proceed with any value. The warning exists because
volume errors are among the most dangerous data-entry mistakes: a user
who accidentally enters their dose in mg when they meant mL (or vice
versa) could see wildly incorrect level predictions.

\subsubsection{Step 6: Injection Site}

``Where did you take it?'' (Skipped for oral routes.) This step presents a
full-screen, rotatable 3D body model (\code{BodyModel3DView.swift},
built on SceneKit) with four snap-angle buttons (Front, Back, Left,
Right) and free-pan drag support. The user taps on the body surface to
mark their injection site. Tapped zones are highlighted with a filled
circle; a zone label appears below the model.

The body surface is divided into clinical zones (ventrogluteal,
dorsogluteal, vastus lateralis, deltoid, etc.) via a hit-testing engine
(\code{ZoneHitTestEngine.swift}) that maps 3D surface coordinates to zone
identifiers. Zones marked as ``avoid'' in Settings are tinted red and
generate a warning on tap. Zones marked as ``preferred'' are tinted
green.

Why a 3D model rather than a 2D diagram or a text list? Because injection
site rotation is spatial. Users need to remember \emph{where} on their
body they last injected---left ventrogluteal or right, upper quadrant or
lower. A flat list of zone names (``Left VG, Right VG, Left DL\dots'')
does not convey the spatial relationship; a 3D model does. The user is
not asked to learn zone names; they are asked to point at their body.

The step is skippable. A ``Skip'' button allows the user to proceed
without marking a site. This matters because some users may not want to
record injection locations, and the app should not block configuration on
a preference.

\subsubsection{Steps 7--9: Anti-Androgens, Body Type, Profile}

These three steps collect supporting information. The anti-androgen step
offers four options (None, Spironolactone, Cyproterone/GnRH agonist,
Finasteride/Dutasteride) with a one-line medical note for each (e.g.,
``Spironolactone blocks androgen receptors and may affect estradiol
clearance''). The body-type step offers three silhouettes for the 3D
model. The profile step collects a name and optional photo.

These steps are placed late in the sequence because they are not required
for the PK engine to function. The app can predict levels without knowing
about spironolactone (the adjustment is optional) or the user's body
type. They are collected at onboarding because they would be unlikely to
be discovered in Settings later, and their absence would silently degrade
the experience.

\subsubsection{Step 10: Review}

The final step displays a summary card of all selections---medication,
ester, dose, concentration, frequency, last dose date, injection site,
anti-androgens, body type, profile name---and a single button: ``Begin
Tracking.'' Tapping it calls \code{commitOnboarding()}, which writes the
configuration to the medication manager, logs the first dose record, sets
the \code{hasCompletedOnboarding} flag, and dismisses the onboarding
cover to reveal the home screen.

The review step is a commitment moment. The user sees everything they
have configured, in one place, before any data is written. There is no
``back'' button from the home screen to onboarding; the configuration is
durable. If the user needs to change something, they do it in Settings,
not by re-running onboarding.

\subsection{First Empty State}
\label{sec:empty_state}

When the home screen first appears after onboarding, the user sees the
brand gradient background, the PK curve (empty, with only the initial
baseline), the level overlay showing the endogenous baseline (e.g.,
``Estimated 100\,pg/mL''), and the glass card stack at the bottom. If
the user skipped dose entry or if the onboarding did not produce a valid
configuration, the glass stack is replaced by an empty state: a pill
icon, the text ``Set up your medication,'' and a prompt to add medication
details.

The empty state is designed to be aspirational, not discouraging. The
graph is visible in the background---the user can see \emph{where} their
data will appear---but the foreground message acknowledges that there is
nothing to show yet. The action to resolve the empty state (record a
dose, configure medication) is always one tap away.

\subsection{First Dose: Celebration and Feedback}
\label{sec:first_dose}

Recording a dose opens the Quick Dose Sheet (\code{QuickDoseSheet.swift}),
a bottom sheet that slides up from the home screen's action bar. The
sheet is organized as a 2-step flow for injectable and topical routes
(dose entry $\to$ site selection), and a single step for oral routes
(dose entry only).

\textbf{Step 1---Dose entry.} The sheet pre-fills today's date and time,
the user's configured dose, and the user's concentration. A mg/mL
segmented control mirrors the one from onboarding. The dose can be
adjusted with a stepper or typed directly. A date/time picker allows
recording a dose for a different time (for catch-up entries).

\textbf{Step 2---Site selection} (non-oral routes only). A compact version
of the 3D body model appears, pre-rotated to the last-used angle. The
user taps to mark the injection site. Below the model, a row of recent
sites shows where the last few injections were placed, with the option to
re-use a recent site in one tap.

\textbf{Completion.} When the dose is saved, the sheet dismisses and,
after a 0.55-second delay (the duration of the settle spring), confetti
particles burst from the center of the screen and the level overlay pulses
with a celebration glow. The confetti consists of 28 spring-animated
particles (circles, rounded rectangles, capsules) in brand colors
(rose, lavender, gold). The particles fan outward with randomized
velocities, decelerate under simulated gravity, and fade out over 2.5
seconds.

Why confetti? Because dose adherence is the most important behavior the
app can encourage, and positive reinforcement at the moment of action is
more effective than reminders before the fact. The confetti is
celebratory but not gamified---there is no streak counter, no points, no
leaderboard. The reward is the moment itself: a brief, beautiful
acknowledgment that the user did the thing they came to the app to do.

After the celebration, the PK curve updates. The single dose appears as a
rising arc on the graph, peaking at the ester's $t_{\max}$ and then
decaying. The level overlay updates to show the current estimated level
at the peak position. For a first-time user, this is the moment the app
transitions from ``I set up a tracker'' to ``it's tracking me.''

\subsection{Daily Use: The Home Screen}
\label{sec:home_screen}

The home screen (\code{HomeView.swift}) is the app's centerpiece. Its
layout is designed for the 95th-percentile interaction: checking your
current level. That interaction requires zero taps, zero scrolling, and
less than two seconds of visual processing.

\textbf{The graph background.} \code{FlatGraphView} renders the PK curve
as a full-bleed cubic-bezier path over a dark gradient background with
rose and purple ambient glows. The curve's color drifts along a
rose-to-lavender ramp keyed to the current level relative to the
therapeutic ceiling: low levels render in deeper rose, high levels in
lighter lavender. This is ambient information---the user absorbs it
peripherally without needing to read numbers. The curve is animated with
an \code{AnimatableVector} that morphs smoothly between data updates,
preventing the jarring jumps that would occur if the path were recomputed
from scratch on each new dose.

Grid lines, Y-axis labels, and peak/trough reference lines are rendered
at low opacity over the curve. The reference lines are drawn from the
user's therapeutic targets if configured, or omitted entirely if not.

\textbf{The level overlay.} \code{LevelOverlayView} is positioned in the
top-left corner, floating over the graph. It displays:
\begin{itemize}[nosep,leftmargin=1.4em]
\item The word ``Estimated'' in small type
\item The current level as a large animated number (e.g., ``284''),
  rendered with \code{AnimatedCounter} for smooth digit transitions
\item The unit (``pg/mL'' or ``ng/dL'')
\item The medication name in smaller type
\item A status pill showing the level's relationship to the user's
  therapeutic targets (e.g., ``Within set trough'')
\end{itemize}

The overlay is the answer to the question ``what is my level right
now?''---the most frequent question the app receives. The design
prioritizes legibility over density: one number, not a dashboard.

\textbf{Scrubbing.} Touching and dragging anywhere on the graph area
activates a scrub mode: the level overlay updates in real time to show
the estimated level at the touched time point, the date appears below the
level value, and a vertical indicator line marks the position. Lifting
the finger returns the overlay to the current-time value. Scrubbing
solves the ``what was my level on Tuesday?'' question without requiring a
separate date-picker interface.

\textbf{The unified glass stack.} \code{UnifiedStackView} is a
glass-morphism card (iOS 26 Liquid Glass with \code{ultraThinMaterial}
fallback) containing stacked rows separated by hairline dividers:

\begin{enumerate}[nosep,leftmargin=1.4em]
\item \textbf{Insight row} (\code{InsightRowView}): Contextual text about
  the current level estimate. At trough: ``Estimated 284\,pg/mL at
  trough---within your set trough range of 100--200\,pg/mL.'' Near peak:
  ``Estimated 612\,pg/mL near peak---above your set peak ceiling of
  600\,pg/mL.'' If a safety flag is active, an additional note follows
  with the clinical association. If spironolactone is configured, a
  clearance note appears: ``Spironolactone may increase estradiol
  clearance by approximately 15\%---this estimate accounts for that
  adjustment.''

\item \textbf{Mood row} (\code{MoodRowView}): Today's journal entry,
  showing the mood emoji, mood name, and a preview of the text. If no
  entry has been recorded, the row reads ``Add mood / Tap to journal''
  and the entire row is tappable to open the journal composer.

\item \textbf{Next dose row} (\code{NextDoseRowView}): The day name of the
  next scheduled dose (``Wednesday''), the number of days remaining, and
  a dot track showing progress through the current cycle. For a 7-day
  cycle on day~3, three filled dots and four unfilled dots are shown. The
  dot track is a spatial representation of cycle position that can be
  read without knowing the numbers.

\item \textbf{Action bar} (\code{ActionBarView}): Two buttons: ``Record
  dose'' (primary, rose-tinted) and ``Journal'' (secondary). These are
  the two actions the home screen is designed to facilitate.
\end{enumerate}

Each row enters with a staggered spring animation (0.07-second delay
between rows, settle spring with 0.55-second duration and zero bounce).
The stagger is fast enough to feel like a single cascade rather than a
sequence of individual animations. On Reduce Motion, all rows appear
immediately.

\textbf{Why four rows?} Each row earns its place by answering a distinct
user question: ``How am I doing?'' (insight), ``How do I feel?''
(mood), ``When's my next dose?'' (next dose), ``What can I do?''
(action bar). Adding more rows would dilute the home screen's primary
function as a two-second check; removing any of the four would leave a
question unanswered. Four is the number that emerged from iteration, not
from a principle; it is the pragmatic answer to ``how much information
fits before the card becomes a dashboard.''

\subsection{Calibration: The First Blood Test}
\label{sec:calibration}

Blood test entry lives in Settings, under History $\to$ Blood Tests. The
composer (\code{BloodTestComposerView.swift}) collects three fields:
date of draw, measured level, and hormone type (estradiol or
testosterone). When saved, the app computes:

\[
\alpha = \frac{C_{\text{measured}}}{C_{\text{predicted}}}
\]

where $C_{\text{predicted}}$ is computed by running the PK model at the
exact draw date \emph{without} calibration applied---the calibration
engine never feeds on its own output, preventing circularity. The factor
is clamped to $[0.2, 5.0]$ to exclude clinically implausible values that
would indicate data-entry error.

The calibration is applied silently. The user does not see $\alpha$
directly unless they view the blood test detail; what they see is that
subsequent level estimates on the home screen are closer to their
measured values. The Labs Chart view
(\code{LabsChartView.swift}) overlays blood test results on the PK
prediction curve, showing the measured values as discrete points and the
model prediction as a continuous line. This is the single most direct
feedback the app provides about its own accuracy: the user can see, at a
glance, how closely the model tracks their actual measurements.

The calibration UX embodies the ``honest uncertainty'' principle
(Section~\ref{sec:uncertainty}). The factor is displayed, the model's
population prediction is visible alongside the calibrated prediction, and
the user can see the magnitude of the adjustment. Nothing is hidden;
nothing is oversold.

\subsection{The Time Tab: Tenure and Constellations}
\label{sec:time_tab}

The Time tab (\code{TimeView.swift}) is the app's most conceptually
ambitious screen. It departs from the home screen's functional density in
favor of ambient, exploratory visualization. Its design philosophy is
that some information about HRT is better \emph{absorbed} than
\emph{read}---that the user's understanding of their own timeline
develops through repeated exposure rather than through a single
explanatory session.

\textbf{Constellation point-clouds.} The upper portion of the screen
contains two large, interactive particle systems: ``Mind''
(brain-shaped point cloud, indigo palette) and ``Body''
(silhouette-shaped point cloud, amber palette). Each cloud is composed of
hundreds of particles rendered in 3D with subtle drift animation. The
clouds are tappable: tapping Mind navigates to mood and cognitive-change
content; tapping Body navigates to physical-change content. A body-system
strip between them summarizes other systems (Skin, Cycle) with smaller
indicator cards.

The cloud density and articulation are driven by the user's months on HRT.
A new user sees sparse, loosely-structured clouds; a veteran sees dense,
clearly articulated shapes. The progression is continuous, not stepped at
tier boundaries---every additional month adds structure. This is ambient
personalization: the user is not told ``your cloud is denser now,'' but
the Time tab feels richer with use, countering the staleness that affects
many health apps' static dashboards.

\textbf{Tenure tiers.} The content model (\code{TimeContentEngine}) is
parameterized by \code{TenureTier}, a 7-level enumeration:
\code{beginning} (0--2\,months), \code{emerging} (2--6), \code{inMotion}
(6--12), \code{settled} (12--24), \code{deepening} (24--60),
\code{established} (60--120), and \code{veteran} (120+). Each tier
unlocks specific content about expected changes, surveillance
recommendations, and body-system timelines. The tier system is documented
in a dedicated design file (\code{TIME\_TIERS\_DESIGN.md}).

Tier transitions trigger a celebration overlay---a brief animated
acknowledgment that the user has reached a new stage. Unlike the confetti
on dose recording, which is immediate and functional, the tier
celebration is reflective: it marks accumulated time, not a single
action.

\textbf{``Up Next'' drawer.} A narrow drawer peeks from the bottom of the
screen, showing only an upward chevron. Pulling it upward triggers a
staggered cascade: the Mind and Body sections lift, shrink, and blur away
(Mind leading, Body following at a greater distance), revealing a list of
upcoming milestones and expected changes. The cascade is the most complex
animation in the app, and it is designed to feel like opening a door: the
content that was ambient becomes focused, and the content that was focused
(the constellations) recedes. On Reduce Motion, the cascade degrades to a
cross-fade.

\subsection{The Journal: Reflection}
\label{sec:journal}

The Journal tab (\code{JournalView.swift}) is the app's reflective
surface. Where the home screen answers ``what is my level'' and the Time
tab answers ``what is happening to my body,'' the Journal answers ``how
do I feel about it.''

\textbf{Composer.} The journal composer
(\code{JournalComposerView.swift}) presents a 7-mood picker
(Happy, Calm, Neutral, Tired, Anxious, Sad, Energetic) rendered as emoji
in a horizontal row, a set of common HRT symptom toggles (skin changes,
hot flashes, fatigue, headache, energy boost, mood swings, breast
tenderness, nausea), and a free-text field. Each entry is timestamped and
automatically tagged with the estimated hormone level at the time of
writing.

The level tagging is automatic and non-optional. The user does not choose
to attach a level reading; the app attaches it because the relationship
between subjective experience and hormone level is precisely what the
Journal is designed to help the user discover. A user who notices that
their ``Anxious'' entries cluster around trough days is seeing a pattern
that the app does not need to point out explicitly.

\textbf{Timeline.} Entries are displayed in a flat timeline grouped by
calendar day, newest first. Each day group shows the date, the primary
mood emoji and name, any symptom tags as colored pills, and a text
preview. Tapping an entry opens a detail sheet. The timeline is designed
for scanning, not searching---its purpose is pattern recognition across
weeks and months, not retrieval of a specific entry.

\subsection{Configuration: Settings}
\label{sec:settings}

The Settings tab (\code{SettingsView.swift}) is organized as a scrollable
list of sections, each with a header and a set of tappable rows or inline
controls. The section order reflects frequency of access: medication
configuration appears near the top; data export and developer settings
near the bottom.

\textbf{Medication management.} Type, method, ester, and dose schedule
can be changed at any time. Changing the medication type or ester
re-triggers the PK model with the new constants; historical dose records
retain their original ester assignments, so a user who switches from oral
to injectable does not see their oral doses reinterpreted with injectable
kinetics.

\textbf{Body zones.} Preferred and avoided injection zones are configured
on the 3D body model, with the same interface used during onboarding. The
zone configuration screens (\code{PreferredZonesEditorView.swift},
\code{AvoidZonesEditorView.swift}) are full-screen editors that allow
multi-select zone marking. The zone preferences feed into the Quick Dose
Sheet's site-selection step, where avoided zones are tinted red and
preferred zones green.

\textbf{Therapeutic targets.} Three numeric fields: trough minimum, trough
maximum, and peak ceiling. These are the only user-configurable clinical
parameters in the app. They draw the reference bands on the home screen
graph and gate the comparative text in the level overlay. The fields are
accompanied by a note: ``Targets are for reference only and do not
constitute medical advice.''

\textbf{History.} Separate list views for dose history
(\code{DoseHistoryView.swift}), blood test history
(\code{BloodTestHistoryView.swift}), and screening history
(\code{ScreeningHistoryView.swift}). Each is a chronological list with
summary cards. The blood test history includes the Labs Chart view,
which overlays measured values on the predicted curve---the most direct
accuracy feedback the app provides.

\textbf{Privacy.} Face ID / passcode lock with configurable auto-lock
timing (immediate, 1\,minute, 5\,minutes, 15\,minutes, 1\,hour). The
lock is implemented as a full-screen overlay (\code{AuthenticationView})
that appears when the app backgrounds, gated by
\code{AuthenticationManager}. Health data is sensitive, and the lock is a
binary control: either the app is locked or it is not. There is no
partial lock that hides some data and shows other data, because partial
locking would create a false sense of security while adding interface
complexity.

% ======================================================================
\section{Encoded Design Decisions}
\label{sec:encoded}

This section documents the design tokens, patterns, and conventions that
ship in the app. It is the UX analog of the PK paper's Encoded Constants
table (Table~2 in~\cite{xtpk}): a reference for what is in the
application, where it came from, and what confidence attaches to it.

\subsection{Color System}
\label{sec:color_system}

\begin{table}[h]
\centering
\small
\begin{tabular}{@{}llll@{}}
\toprule
\textbf{Token} & \textbf{Light} & \textbf{Dark} & \textbf{Role} \\
\midrule
\code{accent}           & \#D4537E & \#E8709A & Primary brand, buttons, links \\
\code{accentSecondary}  & \#534AB7 & \#7E74D8 & Secondary brand, time-mind \\
\code{successBrand}     & \#3B8C6E & \#64DC9E & Positive indicators \\
\code{glassTint}        & \#F0ECF8 & \#0A0814 & Glass card background tint \\
\code{levelLow}         & \#99345A & \#C2557E & Graph curve at low levels \\
\code{levelHigh}        & \#8875D4 & \#B0A4EA & Graph curve at high levels \\
\midrule
\multicolumn{4}{@{}l@{}}{\textbf{Time category palettes (base / surface / border / text)}} \\
Mind    & \#534AB7 / \#EEEDFE / \#CECBF6 / \#3C3489 & \#7E74D8 / \#1E1C3A / \#3D3880 / \#ABA7F0 & Cognitive content \\
Skin    & \#0F6E56 / \#E1F5EE / \#9FE1CB / \#0A4F3E & \#2DB88E / \#0C2420 / \#1A5C46 / \#7FE0C4 & Integumentary content \\
Body    & \#854F0B / \#FAEEDA / \#FAC775 / \#5C360A & \#D4862B / \#2E1C05 / \#6B3D08 / \#EFB25F & Somatic content \\
\bottomrule
\end{tabular}
\caption{Core color tokens. Every color in the application is defined via
\code{Color(lightHex:, darkHex:)}, with explicit light and dark variants.
No color relies on automatic inversion.}
\label{tab:colors}
\end{table}

The color system's design constraint is that every color must work on
both light and dark backgrounds, since the app uses a dark brand gradient
for onboarding and the home screen background but standard system
backgrounds for Settings and sheets. The solution is explicit pairing:
every color is authored twice, once for each appearance. The cost is
maintenance (adding a color means adding two hex values); the benefit is
that no screen looks like an afterthought in either mode.

\textbf{Why rose and lavender?} The brand gradient (rose \#D4537E to
lavender \#8875D4) was chosen to be distinctive without being gendered.
The colors reference medical association (red/pink $\to$ health,
purple $\to$ pharmacy) without mapping to binary gender coding (blue for
testosterone, pink for estrogen). The gradient appears on the onboarding
flow and as an ambient brand marker, never as a functional color.

\subsection{Motion System}
\label{sec:motion}

\begin{table}[h]
\centering
\small
\begin{tabular}{@{}llrrl@{}}
\toprule
\textbf{Preset} & \textbf{API} & \textbf{Duration} & \textbf{Bounce} & \textbf{Use} \\
\midrule
\code{settle} & \code{spring(duration:bounce:)} & 0.55 & 0.00 & Entrance animations, card reveals \\
\code{pop}    & \code{spring(duration:bounce:)} & 0.40 & 0.15 & Interactive responses \\
\code{gentle} & \code{spring(duration:bounce:)} & 0.50 & 0.08 & Warm transitions \\
\code{tap}    & \code{spring(duration:bounce:)} & 0.25 & 0.05 & Button press feedback \\
\code{burst}  & \code{spring(duration:bounce:)} & 0.60 & 0.28 & Celebrations, confetti \\
\bottomrule
\end{tabular}
\caption{Named spring presets. All use the WWDC23 \code{spring(duration:bounce:)}
API. The \code{settle} spring (0.55\,s, 0 bounce) is the default entrance
curve; its zero-bounce decay feels natural without calling attention to
itself.}
\label{tab:springs}
\end{table}

\textbf{Entrance animations.} The \code{EntranceModifier} fades and lifts
views into place from below (16pt offset) with a configurable delay.
Stacked cards use per-index staggering at 0.07-second intervals via
\code{.entrance(index:)}, producing a cascade that reads as a single
motion rather than a sequence. The stagger interval was tuned by feel: at
0.05\,s the cascade is imperceptible; at 0.10\,s it feels slow. 0.07\,s
sits at the threshold where the cascade is noticeable but not
distracting.

All entrance animations are gated on
\code{@Environment(\.accessibilityReduceMotion)}. When Reduce Motion is
enabled, views appear immediately at their final position.

\textbf{Tactile buttons.} \code{PressableButtonStyle} applies a 0.95
scale reduction and 0.92 opacity reduction on press, animated with the
\code{tap} spring (0.25\,s, 0.05 bounce), and fires a soft-impact haptic
(\code{UIImpactFeedbackGenerator(style: .soft)}). The combination of
visual compression and haptic feedback makes buttons feel physical---a
property that matters for the Record Dose button, which is the most
frequently tapped control in the app.

\textbf{Celebration effects.} Two effects fire on dose recording: a
confetti burst of 28 particles (circles, rounded rectangles, capsules in
brand colors) with randomized trajectories, gravity, and fade; and a
celebration glow (\code{CelebrationGlowModifier}) that pulses the level
overlay with a PhaseAnimator-driven bloom. The confetti is gated on
Reduce Motion; the glow degrades to a simple opacity pulse.

\subsection{Typography}
\label{sec:typography}

\begin{table}[h]
\centering
\small
\begin{tabular}{@{}lll@{}}
\toprule
\textbf{Token} & \textbf{Base Style} & \textbf{Role} \\
\midrule
\code{statValue}  & Large title, rounded design & Home screen level readout \\
\code{hero}       & Title, light weight        & Onboarding welcome \\
\code{readout}    & Title 2, rounded design     & Secondary numeric displays \\
\code{eyebrow}    & Caption 1, uppercased       & Section labels \\
\code{pill}       & Caption 2, medium weight    & Status badges \\
\code{detail}     & Subheadline                 & Form field labels \\
\code{label}      & Body                        & Primary body text \\
\code{numeric}    & Title 3, monospaced         & Dose entry fields \\
\bottomrule
\end{tabular}
\caption{Role-based font tokens. Every token is built on a Dynamic Type
text style, so all text scales with the user's system font size
preference. Custom tokens add weight and design adjustments on top of the
base style.}
\label{tab:typography}
\end{table}

The typography system's constraint is that fonts must be legible at very
large sizes (the home screen level readout is the largest text in the
app, and it must work at Accessibility sizes) and at very small sizes
(status pills, axis labels). The solution is role-based tokens that
abstract away the underlying font metrics. If a token needs adjustment, it
is changed in one place; all views that use the token inherit the change.

\subsection{Surface Language}
\label{sec:surfaces}

\textbf{Glass morphism.} The unified card stack on the home screen uses a
glass effect: on iOS 26, \code{.glassEffect} produces the system Liquid
Glass material; on earlier iOS versions, \code{ultraThinMaterial} with a
tinted background provides a visual fallback. The glass card carries
rounded corners (\code{Theme.Radius.lg}) and a subtle border
(\code{Color.white.opacity(0.2)}). The effect is decorative, not
functional---the card would work identically on a solid background---but
the glass aesthetic serves a purpose: it visually separates the card's
information (read, actionable) from the graph behind it (ambient,
decorative), reinforcing the home screen's information hierarchy.

\textbf{Time cards.} Content in the Time tab uses opaque pastel cards
colored with the category palette's surface color. These cards are
tappable navigation targets, not ambient displays, so they use solid
backgrounds rather than glass. The distinction is consistent: glass means
``this is an overlay on ambient content''; solid means ``this is a
discrete piece of information.''

\textbf{Dividers.} Separators between rows in the unified stack are
0.5pt rectangles at 35\% opacity of the border color. Thinner than the
system default (\code{Divider()}), they mark boundaries without
dominating the card. The half-point weight is below the threshold where a
line reads as a structural element; it reads as a seam.

\subsection{Presentation Rules}
\label{sec:presentation_rules}

The file \code{HormoneLevelPresentation.swift} (approximately 300 lines)
is the single point where hormone level data becomes user-facing text. It
exports:

\begin{itemize}[nosep,leftmargin=1.4em]
\item \code{LevelStatus} enum: 6 cases from \code{criticallyLow} to
  \code{high}
\item \code{SafetyFlag} enum: 4 clinical threshold crossings
\item \code{ModelDeviation} enum: how calibrated prediction compares to
  population prediction
\item \code{RangeSource} enum: \code{userSet} or \code{modelDerived}
\item Formatting functions that accept these enums and return
  localization-ready strings
\end{itemize}

The design rule is that no other file in the application may contain a
hardcoded string that describes a hormone level in comparative or
clinical terms. All such strings flow through this file. This
centralization exists for three reasons: it keeps the presentation rules
auditable (every instance of ``below target'' can be found in one file);
it ensures consistency (two screens describing the same level status will
use the same phrasing); and it simplifies localization (the
presentation-rules file is the natural target for translation).

\subsection{Navigation Architecture}
\label{sec:navigation}

\begin{table}[h]
\centering
\small
\begin{tabular}{@{}llll@{}}
\toprule
\textbf{Tab} & \textbf{Icon} & \textbf{Function} & \textbf{Interaction Model} \\
\midrule
Home    & house          & Check level, record dose & Glance + action \\
Journal & book           & Record and review moods  & Compose + scan \\
Time    & wand.and.stars & Explore timeline         & Browse + discover \\
Settings& gearshape      & Configure                & Edit + review \\
\bottomrule
\end{tabular}
\caption{The 4-tab architecture. Each tab has a distinct interaction
model: Home is for quick actions, Journal for composition and review,
Time for browsing, Settings for editing. The models are intentionally
different---a tab that duplicated Home's glance-and-act pattern would be
redundant.}
\label{tab:tabs}
\end{table}

The 4-tab structure is a standard iOS pattern (\code{TabView} with
\code{.automatic} style). The tabs are not customizable; their order is
fixed. The decision to use four tabs rather than three or five was driven
by content, not by principle: Home (dashboard), Journal (reflection), and
Settings (configuration) were the three non-negotiable screens, and Time
(exploration) earned its place as the fourth by covering a user need---``what's
happening to me over time''---that none of the other three addressed.

% ======================================================================
\section{Validation: Persona Journeys}
\label{sec:personas}

Section-level walkthroughs of screens describe the interface as designed.
Persona journeys describe the interface as \emph{used}---by specific
people, on specific regimens, encountering specific situations. This
section adopts the persona-validation approach of the companion PK
paper (Section~4.4 of~\cite{xtpk}), exercising UX edge cases rather
than PK edge cases.

\begin{longtable}{@{}p{2.6cm}p{4.5cm}p{7.4cm}@{}}
\toprule
\textbf{Persona} & \textbf{Scenario} & \textbf{What it tests, and the result} \\
\midrule
\endhead
Maya, 24 & 4\,mg EV IM q7d, first-time user, completes full onboarding
 & \textbf{Onboarding completeness.} Maya sees 11 screens (injectable
   route), each presenting one decision. The ester picker shows EV's
   half-life (2.6\,d) and peak ($\sim$2\,days). She selects
   40\,mg/mL concentration (marked ``most common''), enters 4\,mg dose
   (0.10\,mL auto-computed), taps her left ventrogluteal zone on the 3D
   model, and reviews her configuration before committing. First dose is
   logged; confetti fires; the PK curve shows a single-dose arc peaking at
   2.1\,d. The level overlay reads ``Estimated 307\,pg/mL'' near peak.
   \good{Pass.} \\[0.4em]

Devon, 31 & 100\,mg TC q7d, switches to 200\,mg q14d at month~6
 & \textbf{Regimen change UX.} Devon opens Settings $\to$ Dose Schedule,
   changes frequency from 7 to 14 days, changes dose from 100 to
   200\,mg. The PK model recomputes: $C_{\mathrm{avg}}$ is preserved,
   trough falls, peak rises, peak/trough widens from $1.74\times$ to
   $3.51\times$. The home screen curve reshapes to show the longer
   interval. The next-dose countdown updates to 14 dots. Devon does not
   need to re-enter onboarding or re-log past doses; the configuration
   change applies forward. \good{Pass.} \\[0.4em]

Priya, 47 & 2\,mg oral E2 BID for 4 months, then switches to 4\,mg EV q7d
 & \textbf{Route switch.} Priya changes her medication type (still
   Estrogen) and method (Injectable) in Settings. The app prompts for
   ester, concentration, and dose---the fields that differ between routes.
   Historical oral doses retain their oral kinetics ($F = 0.05$, fast
   rates); new injectable doses use EV kinetics. The graph shows a
   continuous curve across the transition. Priya does not need to delete
   her history to change routes. \good{Pass.} \\[0.4em]

Sam, 19 & 50\,mg/day T gel, misses $\sim$2 days/week
 & \textbf{Missed-dose handling.} Sam's dose record status is ``missed''
   for the skipped days. The PK model computes the curve without those
   contributions. $C_{\mathrm{avg}}$ drops to 637\,ng/dL; on missed days,
   the level overlay shows a lower value (roughly 20\% below average,
   63\% below normal trough). The next-dose row continues to show
   ``Today'' since gel is daily; the dot track shows one dot (daily
   cycle). Sam sees the consequence of missed doses without the app
   moralizing about adherence. \good{Pass.} \\[0.4em]

Rowan, 38 & 4\,mg/day compounded E2 cream, custom concentration
 & \textbf{Custom concentration entry.} Rowan selects ``Custom'' in the
   concentration step, types their compounded concentration (e.g.,
   2.5\,mg/mL), and proceeds. The dose step computes mL from the custom
   concentration correctly. The route pseudo-ester system resolves
   ``Cream'' to the transdermal estradiol constants. The graph uses the
   24-hourly micro-dose model (Equation~4 in~\cite{xtpk}). The custom
   concentration value persists and appears in Settings.
   \good{Pass.} \\[0.4em]

Kai, 52 & 237\,mg oral TU BID, 10-day interruption (travel, ran out)
 & \textbf{Washout and re-accumulation.} Kai records no doses for
   10~days. The PK curve falls to the endogenous floor within $\sim$1~day
   (oral TU's fast kinetics: $k_e = 8.305$\,day$^{-1}$, $t_{1/2} =
   2.0$\,h). No phantom drug survives the gap. When Kai resumes, the
   first new dose appears on the curve at the correct date; the
   re-accumulation to steady state takes $\sim$2\,days. The level overlay
   reads near zero during the gap, which is clinically accurate. Kai does
   not need to tell the app about the interruption; the absence of dose
   records is the signal. \good{Pass.} \\[0.4em]

Jules, 28 & 100\,mg TE IM q7d, first blood test at month~3
 & \textbf{Calibration UX.} Jules enters a blood test result:
   680\,ng/dL at trough. The app computes $\alpha = 1.20$ (measured
   20\% above population prediction). The Labs Chart shows the measured
   point above the predicted curve. Subsequent level estimates are scaled
   by 1.20; the home screen now reads $\sim$680 at trough rather than the
   population prediction of $\sim$566. The $\alpha$ factor is visible in
   the blood test detail but not in the home screen overlay. Jules does
   not need to understand $\alpha$ to benefit from it. \good{Pass.}
   \\[0.4em]

Nadia, 35 & 5\,mg EC q10d---an interval nobody studied
 & \textbf{Extrapolation outside the calibration envelope.} Nadia
   configures a custom 10-day frequency. The PK model handles it---the
   superposition equation (Equation~2 in~\cite{xtpk}) is agnostic to
   interval---but the resulting curve sits outside the range where the
   EC constants were anchored (steady-state q7d). The trough at q10d is
   $\sim$40\% below the q7d trough for the same dose. The app does not
   warn Nadia that her interval is unconventional; it computes the
   prediction and presents it. The companion PK paper flags this as
   ``structurally sound but inherits every assumption of the fit'';
   the interface does not surface that caveat. \good{Pass, with caveat.}
   \\[0.4em]
\bottomrule
\end{longtable}

All eight personas complete their journeys without interface friction.
The personas test properties that section-level screen descriptions
cannot: the discoverability of settings for regimen changes (Devon), the
handling of route switches with mixed-dose histories (Priya), the
consequence visualization of missed doses (Sam), custom-value entry paths
(Rowan), washout handling during gaps (Kai), the transparency of
calibration (Jules), and the behavior of the interface outside its design
envelope (Nadia).

% ======================================================================
\section{Discussion: Design Trade-offs and Limitations}
\label{sec:discussion}

This section describes trade-offs that were consciously accepted and
limitations that remain unaddressed. It is the UX analog of the companion
PK paper's Discussion (Section~5 of~\cite{xtpk}), which quantifies
structural limitations rather than hiding them.

\subsection{Onboarding Length vs.\ Setup Completeness}
\label{sec:onboarding_length}

The 11-step onboarding is long relative to consumer app norms. Many
health apps complete setup in 3--5 screens, accepting that additional
configuration will happen later in Settings. Isocline's position is that later
configuration does not happen---users who skip a setting at onboarding
rarely discover it in Settings, and the app's predictions silently degrade
as a result.

The trade is that onboarding takes 3--5~minutes, which is above the
median for consumer app setup, in exchange for the app being fully
configured at first launch. The cost is potential drop-off during
onboarding; the benefit is that users who complete it have a correctly
parameterized PK model from day one. Whether the drop-off rate justifies
the completeness is an empirical question the current design cannot
answer without analytics the app does not collect.

\subsection{The Graph vs.\ The Number}
\label{sec:graph_vs_number}

The home screen privileges the level overlay (one number) over the PK
curve (the full time series). This is a deliberate choice: the number
answers the most frequent question immediately; the curve answers
follow-up questions on demand (via scrubbing). The cost is that users may
not discover the curve's information density---the peak/trough pattern,
the accumulation across doses, the shape of the decay---because the
interface does not push them to look at it.

An alternative design would foreground the curve and derive the number
from it, making the time-series the primary visual and the current-value
readout secondary. This would better serve users who think in terms of
curves and cycles, at the cost of slower legibility for the
glance-and-go interaction. The current design favors the novice; a
future iteration might introduce a mode that favors the experienced user.

\subsection{Medical Information Without Medical Advice}
\label{sec:medical_info}

Section~\ref{sec:state_facts} describes the principle. This section
describes its cost. The app's refusal to advise means that a user who
\emph{wants} directive guidance---``your levels are too high, skip your
next dose''---will find Isocline unsatisfying. The app will tell them their
level, tell them it exceeds a clinical threshold, and state the
association. It will not tell them what to do about it.

This is the correct design for a consumer health application operating
without clinical oversight, but it creates a gap between what some users
expect from a ``health app'' (directive, prescriptive) and what Isocline
provides (informative, associative). Bridging that gap without crossing
the line into advice is an unsolved design problem; the current interface
acknowledges it by erring on the side of reticence.

\subsection{Tenure Tiers as Content Gating}
\label{sec:tier_gating}

The Time tab's tenure-tier system withholds information from early-stage
users. A user in month~1 does not see year-3 surveillance content; a user
in year~5 does not see month-1 ``first changes'' content. The gating is
intentional: information presented before it is relevant is noise.

But the system assumes a standard progression through HRT that may not
match every user's experience. A user who starts, stops, and restarts
HRT---or who is on a non-standard protocol---may see content gated behind
a tier that does not reflect their clinical reality. The tenure counter
(\code{activeMonthsOnHRT}) counts months since the first dose record; it
does not account for gaps. A user who took HRT for 6~months, stopped for
2~years, and restarted would be in the \code{inMotion} tier (6--12~months)
despite having only a few weeks of current experience. This is a known
limitation of the tier model; addressing it would require a more
sophisticated treatment-interruption model that the current design does
not include.

\subsection{Calibration as a Single Scalar}
\label{sec:calibration_scalar}

The per-user calibration factor $\alpha$ (Section~\ref{sec:calibration})
is the UX equivalent of the PK paper's ``pragmatic IIV analog''
(Section~5.5 of~\cite{xtpk}). It is a single scalar multiplier on the
entire model output. It cannot represent parameter-specific variability: a
user with fast clearance but normal volume of distribution, or a user
with high bioavailability but slow absorption, both get the same
one-number adjustment.

The UX consequence is that calibration improves accuracy on average but
may not improve it at every point in the cycle. A user whose $\alpha$ was
derived from a trough measurement may find that their peak predictions
are \emph{less} accurate after calibration, because the factor that
corrects the trough over-corrects the peak. The interface does not
surface this possibility; it shows the calibrated curve as a single line,
not as a band that widens to reflect parameter uncertainty. This is a
conscious simplification: showing uncertainty bands would be more honest
but also more confusing for users who are not accustomed to interpreting
confidence intervals.

\subsection{What Remains Unaddressed}
\label{sec:unaddressed}

The following limitations are structural---they cannot be addressed by
refining the existing interface and would require new features or
architectural changes.

\textbf{No provider sharing.} The app has no mechanism for sharing data
with a healthcare provider. A user who wants to show their PK curve to
their endocrinologist must take a screenshot. A provider-facing view or
export format (PDF report, FHIR integration) would close this gap but
introduces privacy and regulatory complexity (HIPAA compliance, data
transmission security) that the current architecture does not address.

\textbf{No multi-medication tracking.} The app tracks one medication at a
time. A user on both estradiol and progesterone, or testosterone and
finasteride, must choose which to track. The PK model could handle
multiple medications in principle, but the interface would need
substantial restructuring: a second level overlay, a second curve or a
combined curve, and a medication switcher. The current single-medication
design is simpler but excludes combination therapy users.

\textbf{No Apple Health integration.} The app does not read from or write
to HealthKit. Blood test results, dose records, and journal entries live
only within Isocline. Integration with Apple Health would allow the app to
pull lab results entered in other apps and push dose records for use by
other health tools. The barrier is not technical but design-level: Apple
Health's data model does not have native types for ``intramuscular
injection of estradiol valerate at 40\,mg/mL concentration,'' so mapping
Isocline's rich dose records to HealthKit's simpler types would lose
information.

\textbf{No social or community features.} The app is intentionally
single-user. There is no mechanism for comparing experiences, sharing
insights, or connecting with other users. This is a product decision, not
an oversight: community features introduce moderation burden, privacy
risk, and the potential for peer advice to substitute for clinical
guidance. But it means the app does not address the social dimension of
HRT---the desire to know that other people are experiencing similar
things---which is a significant part of the user need.

\textbf{Limited gender model.} The app's medication picker offers two
paths: estrogen-based (feminizing) and testosterone-based
(masculinizing). This binary framing reflects the current clinical
landscape of gender-affirming HRT but does not accommodate non-binary
users who may be on regimens that do not map cleanly to either path, or
who may prefer not to characterize their HRT in gendered terms. The
clinical constants are correct regardless of the label, but the
onboarding language could be more inclusive.

\textbf{English-only.} The app is monolingual. The presentation-rules
file (\code{HormoneLevelPresentation.swift}) was designed for
localization---classifications are enums, display strings are
centralized---and a localization pass is in progress on the current
development branch, but as of this writing, every user-facing string is
English. For an app targeting a global user base, this is the most
significant accessibility gap.

\textbf{Graph accessibility.} The PK curve is not navigable by VoiceOver.
A VoiceOver user can hear the current level estimate (the overlay is a
standard accessibility element) but cannot explore the curve by scrubbing
or hear a description of its shape. Making a continuous mathematical
curve accessible to screen-reader navigation is an unsolved problem in
iOS accessibility; the app inherits that limitation.

\subsection{Design Decision Confidence Hierarchy}
\label{sec:confidence}

Analogous to the PK paper's source quality hierarchy (Section~5.3
of~\cite{xtpk}), design decisions in Isocline rest on a hierarchy of
confidence:

\begin{enumerate}[nosep,leftmargin=1.4em]
\item \textbf{Platform-standard} (highest confidence). Decisions that
  follow iOS Human Interface Guidelines: tab bar navigation, modal sheet
  presentation, standard gesture recognizers, Dynamic Type, Reduce
  Motion, VoiceOver labeling. These patterns are tested across millions
  of apps and billions of user sessions.
\item \textbf{Clinically-grounded.} Decisions that follow from the PK
  model or from published clinical data: the choice of esters, the
  therapeutic range defaults, the safety flag thresholds, the tenure tier
  boundaries. These are anchored to the same sources as the PK paper's
  constants.
\item \textbf{Reasoned but untested.} Decisions based on design reasoning
  without formal user testing: the 11-step onboarding length, the 4-row
  card stack composition, the mg/mL toggle design, the confetti
  celebration. These are internally consistent and follow from the design
  principles, but their effectiveness has not been measured.
\item \textbf{Placeholder} (lowest confidence). Decisions known to be
  suboptimal but retained because a better alternative is not yet
  available: TP's uncalibrated $t_{\max}$ (not presented to users, but
  the curve shape is wrong for the first $\sim$24\,hours), the single-$\alpha$
  calibration (correct on average, not at every cycle point), the lack
  of treatment-interruption modeling for tenure tiers.
\end{enumerate}

\subsection{Limitation Surfacing}
\label{sec:limitation_surfacing}

A recurring question in the design of Isocline is: \emph{which} of the app's
known limitations should be surfaced to the user, and \emph{how}? The PK
paper documents structural limitations in detail (Sections~5.1--5.5
of~\cite{xtpk}), but the consumer interface does not present these
caveats alongside predictions. The decision not to surface every
limitation is itself a UX choice, and it is worth explaining.

The principle is: surface limitations that the user can act on; document
limitations that they cannot. The calibration factor $\alpha$ is surfaced
because the user provided the blood test and can see the adjustment.
Therapeutic target bands are surfaced because the user configured them and
can change them. Safety flags are surfaced because they describe
actionable clinical associations.

The PK model's structural limitations---TU peak/trough ratio,
transdermal ripple, oral testosterone peak under-prediction, TP
uncalibration---are \emph{not} surfaced because the user cannot do
anything with that information. Telling a testosterone undecanoate user
that ``your peak/trough ratio may differ from the model's prediction due
to a conflict between published sources'' is technically honest but
practically useless: it adds cognitive burden without enabling action.

This is a judgment call, and reasonable designers could disagree. The
counterargument is that transparency about model limitations builds trust
and allows users to calibrate their own confidence in the predictions.
The current design accepts the trade: it prioritizes a clean interface
over exhaustive caveats, on the reasoning that the user who wants to know
about PK model limitations can read the companion paper, and the user who
does not should not have caveats forced into their daily experience.

% ======================================================================
\section{Conclusion}
\label{sec:conclusion}

We have presented the user experience architecture of Isocline, a mobile
application for hormone replacement therapy monitoring. The interface
rests on three principles---progressive disclosure, honest presentation,
and per-user personalization---that operate at every level of the
application, from the 11-step onboarding wizard to the ambient
constellation visualizations on the Time tab.

The UX architecture's key structural features are:

\begin{itemize}[nosep,leftmargin=1.4em]
\item A \textbf{progressive disclosure framework} operating on three
  timescales simultaneously: onboarding (minutes), per-tab depth
  (seconds), and tenure-tier content gating (months to years).
\item A \textbf{presentation-rules file}
  (\code{HormoneLevelPresentation.swift}) that centralizes every
  instance of clinical language, ensuring that the app states facts
  rather than giving advice, gates comparisons behind user-set targets,
  and keeps semantics separate from display for localization.
\item A \textbf{design-token system} (colors, motion, typography,
  surfaces) that encodes clinical semantics as visual properties: the
  graph curve's color ramp tracks the therapeutic range, the Time tab's
  palette system maps body systems to colors, and the motion presets
  encode interaction semantics (settle for arrivals, burst for
  celebrations).
\item A \textbf{per-user calibration interface} that translates a blood
  test result into a single multiplicative factor $\alpha$, absorbing
  the dominant source of individual-level prediction error without
  requiring the user to understand population pharmacokinetics.
\item An \textbf{accessibility architecture} in which Reduce Motion,
  Dynamic Type, and dark/light adaptivity are first principles rather
  than retrofits, with every animation, font, and color gated on the
  appropriate system setting.
\end{itemize}

The eight persona journeys (Section~\ref{sec:personas}) exercise the
architecture across onboarding, regimen changes, route switches, missed
doses, custom configurations, treatment gaps, calibration, and
extrapolation beyond the design envelope. All eight complete without
interface friction.

Design trade-offs are named rather than hidden: the 11-step onboarding is
long but complete; the home screen privileges a single number over the
full curve; the calibration factor is a single scalar that cannot
represent parameter-specific variability; the tenure-tier system assumes
continuous treatment; and the model's structural limitations are not
surfaced in the consumer interface, on the reasoning that users cannot
act on them.

What remains unaddressed---provider sharing, multi-medication tracking,
Apple Health integration, community features, a more inclusive gender
model, localization, and graph accessibility---is stated plainly.

The most important property of a health-tracking interface, in our view,
is the same as the most important property of the pharmacokinetic model
that drives it: \textbf{honesty about what it can and cannot say.} Isocline's
interface is designed to be honest. It reports estimates as estimates,
not as facts. It compares only against user-set targets, not against its
own population defaults. It states clinical associations without
prescribing actions. It names its limitations in documentation, even when
it does not surface them in the interface. This honesty---the explicit
marking of what is designed, what is reasoned, and what is
placeholder---is not the most visible property of the interface, but it
is the one that matters most for an application whose users may make
medical decisions based on its output.

\vfill
\noindent\rule{\linewidth}{0.4pt}\\[0.3em]
{\footnotesize This paper describes a consumer tracking application's
user experience design. It is not clinical guidance, and the interface
design decisions described here should not be interpreted as medical
recommendations. The application's predictions are population estimates
that should not be used to make dosing decisions without laboratory
confirmation.}

% ======================================================================
\begin{thebibliography}{99}

\bibitem{xtpk}
Cortes S.\ (2026). ``Superposed Bateman Biexponentials with
Clearance-Constrained Scaling and Per-User Calibration: A Pharmacokinetic
Framework for Hormone Replacement Therapy Monitoring in the Isocline Mobile
Application.'' August 2026.

\bibitem{hembree2017}
Hembree W.C., et al.\ (2017). ``Endocrine Treatment of Gender-Dysphoric
/Gender-Incongruent Persons: An Endocrine Society Clinical Practice
Guideline.'' \emph{J.\ Clin.\ Endocrinol.\ Metab.} 102(11):3869--3903.

\bibitem{ichihara2020}
Ichihara K., Masumori N., Fujii S., Toda T.\ (2020). ``Pharmacokinetics
of Testosterone Enanthate After Intramuscular Injection for Transgender
Men.'' \emph{Androgens: Clinical Research and Therapeutics}.
\texttt{doi:10.1089/andro.2020.0002}.

\bibitem{bi2018}
Bi Y., et al.\ (2018). ``Population Pharmacokinetic/Pharmacodynamic
Modeling of Depot Testosterone Cypionate in Healthy Male Subjects.''
\emph{CPT: Pharmacometrics \& Systems Pharmacology}.
\texttt{doi:10.1002/psp4.12287}.

\end{thebibliography}

\end{document}
